% =====================================================================
%  main.tex -- short note.
%  Everything lives in this one file: no sections/ directory, no
%  separate preamble, no table of contents. Use this for manuscripts up
%  to roughly fifteen pages whose argument does not divide into
%  independent parts.
%  Build with:  pdflatex main; bibtex main; pdflatex main; pdflatex main
% =====================================================================
\documentclass[11pt]{article}

\usepackage[T1]{fontenc}
\usepackage[utf8]{inputenc}
\usepackage{lmodern}
\usepackage[margin=1in]{geometry}
\usepackage{amsmath,amssymb,amsthm,mathtools}
\usepackage{enumitem,booktabs,array}
\usepackage{microtype}
\usepackage{xcolor}
\definecolor{linkblue}{RGB}{25,65,105}
\usepackage[colorlinks=true,linkcolor=linkblue,citecolor=linkblue,
            urlcolor=linkblue]{hyperref}
\usepackage[capitalise,noabbrev,nameinlink]{cleveref}

\numberwithin{equation}{section}

\newtheorem{theorem}{Theorem}[section]
\newtheorem{lemma}[theorem]{Lemma}
\newtheorem{proposition}[theorem]{Proposition}
\newtheorem{corollary}[theorem]{Corollary}
\theoremstyle{definition}
\newtheorem{definition}[theorem]{Definition}
\theoremstyle{remark}
\newtheorem{remark}[theorem]{Remark}

\newcommand{\R}{\mathbb{R}}
\newcommand{\Z}{\mathbb{Z}}
\newcommand{\dd}{\mathrm{d}}
\newcommand{\norm}[1]{\left\lVert #1 \right\rVert}

\title{A Short Note Template}
\author{[OI]}
\date{\today}

\hypersetup{pdftitle={A Short Note Template},pdfauthor={[OI]}}

\begin{document}

\maketitle

\begin{abstract}
A short note states one result and proves it. There is no table of
contents, no separate preamble file, and no appendix unless a single
deferred computation genuinely needs one. Replace this text with one
self-contained paragraph giving the problem, the result, and its
significance.
\end{abstract}

\section{Introduction}
\label{sec:introduction}

State the problem and the result. In a note, the literature discussion
is one or two sentences, not a survey: name the closest prior result
\cite{example-classical} and say what is new here.

\begin{theorem}[Main theorem]
\label{thm:main}
Let $X$ be a compact metric space and let $\mu$ be a Borel probability
measure on $X$. For every $f \in L^2(\mu)$,
\begin{equation}
\label{eq:main}
  \int_X f \dd\mu \le \norm{f}_{L^2(\mu)} .
\end{equation}
\end{theorem}

\begin{proof}
Apply Cauchy--Schwarz to the pairing of $f$ with the constant function
$1$ and use $\mu(X) = 1$.
\end{proof}

\begin{remark}
\label{rem:sharp}
The constant in \cref{eq:main} is attained by $f = \mathbf{1}_X$.
\end{remark}

\section{Consequences}
\label{sec:consequences}

A note may have a short second section for immediate corollaries, but
it should not accumulate machinery. If a second section needs its own
definitions and standing conventions, the material belongs in the
standard article template instead.

\bibliographystyle{plain}
\bibliography{references}

\end{document}
